\subsection{Reduction to $\indcca$ security of $\KEM$ and PRF security in the standard model}

\todo[inline]{Literally the standard model proof of X-Wing relying on the
    IND-CCA security of the KEM and the KDF being a secure PRF}

\begin{theorem} \label{theorem:kem-sm}
    Let $\nominalgroup = \nominalgroupexpanded$ be a nominal group and \kem be a key encapsulation mechanism, let $H$ be a secure PRF with an output size of $n$ when keyed on $k_1$, and let $\adversary{A}$ be an adversary against the \cca security of \hqb in Definition \ref{def:hqb}.
    Then, there exist $\adversary{B}$,  $\adversary{C}$, $\adversary{D}$, and $\adversary{E}$, such that,
    \[
        \advantage{\cca, \adversary{A}}{\hqb} \leq \advantage{\cca, \adversary{B}}{\kem} + \advantage{\cca, \adversary{C}}{\kem} + \advantage{\prf, \adversary{D}}{H} + \advantage{\prf, \adversary{E}}{H} + 2\delta \, ,
    \]
    where $\delta$ is the correctness bound for \kem.
    The run-times of $\adversary{B}$,  $\adversary{C}$, $\adversary{D}$, and $\adversary{E}$ are roughly the same as of $\adversary{A}$.
    $\adversary{B}$ and $\adversary{C}$ perform at most $q_d$ queries to their own decapsulation oracles.
\end{theorem}

\begin{proof}
    This proof follows closely the proof for \cite[Theorem 1]{PKC:GiaHeuPoe18}.

    \begin{gamefigure}{Games $G_0$ to $G_4$.}{fig:kem-sm-reduction}
        \begin{games}{2}
            \begin{gameenv}{\game $G_0 / \textcolor{\gone}{G_1} / \textcolor{\gtwo}{G_2} / \textcolor{\gthree}{G_3} / \textcolor{\gfour}{G_4}$}
                \State $\sk_1, \pk_1 \sample \kem.\keygen()$
                \State $\sk_2 \sample \honestexponent$
                \State $\pk_2 \assign \nexp(\ngen, \sk_2)$
                \State $\pk \assign (\pk_1, \pk_2)$
                \State $\sk \assign (\sk_1, \sk_2)$
                \State $k^*_1, c^*_1 \sample \kem.\enc(\pk_1)$
                \State $\sk_e \sample \exponent$
                \State $c^*_2 \assign \nexp(\ngen, \sk_e)$
                \State $k^*_2 \assign \nexp(\pk_2, \sk_e)$
                \BeginBox[draw=\gone]
                \State $k_1^* \sample \mathcal{K}$
                \Comment{$G_1 - G_3$}
                \EndBox
                \State $k^* \assign H(\hashinputchallenges)$
                \BeginBox[draw=\gtwo]
                \State  $k^* \sample \bits^{n}$
                \Comment{$G_2 - G_4$}
                \EndBox
                \State $c^* \assign (c_1^*, c_2^*)$
                \State $b' \sample \adversary{A}^{\dec(\cdot)}(\pk, c^*, k^*)$
                \State \Return $b'$
            \end{gameenv}

            \gamebreak

            \begin{gameenv}{\oracledec}
                \State $\ifthen{c = c^*}$
                \State $\quad$ \Return $\bot$
                \State $(c_1, c_2) \assign c$
                \State $k_1 \assign \kem.\dec(c_1, \sk_1)$
                \State $k_2 \assign \nexp(c_2, \sk_2)$
                \State $\ifthen{k_1 \compare \bot}$
                \State \quad \Return $\bot$
                \BeginBox[draw=\gone]
                \State $\ifthen{c_1 = c^*_1}$
                \Comment{$G_1 - G_3$}
                \State $\quad k_1 \assign k_1^*$
                \EndBox
                \State $k \assign H(\hashinput)$
                \BeginBox[draw=\gtwo]
                \State $\ifthen{c_1 = c^*_1}$
                \Comment{$G_2$}
                \State $\quad k \sample \bits^{n}$
                \EndBox
                \State \Return $k$
            \end{gameenv}
        \end{games}
    \end{gamefigure}

    \begin{itemize}
    \item[\textbf{Game 0}]
        This is the standard IND-CCA security game for \hqb and so,

        \[ \prone{\cca^0_{\hqb, \adversary{A}}} = \prone{G_0}. \]

    \end{itemize}

    \begin{itemize}
    \item[\textbf{Game 1}]
        The challenge shared key produced by \kem is replaced with a random key.
    \item[\emph{Claim 1}]
        This change should not be noticeable by the adversary $\adversary{A}$.
        If it is, then we can construct an efficient adversary $\adversary{B}$ against the IND-CCA security of \kem with roughly the same running time as $\adversary{A}$, such that,

        \[
            \abs{\prone{G^{\adversary{A}}_0} - \prone{G^{\adversary{A}}_1}} \leq \advantage{\kem}{\cca, \adversary{B}} + \delta.
        \]

    \item[\emph{Proof}]
        We construct $\adversary{B}$ as in Figure \ref{fig:adversary-B}.
        The adversary $\adversary{B}$ simulates the environment of $\adversary{A}$ by calculating the nominal group components of \hqb itself and embedding the \kem challenge into the \hqb challenge.
        The decapsulation oracle is simulated by $\adversary{B}$ as follows: it can calculate the nominal group part itself and query its own \kem decapsulation oracle on all ciphertexts except $c_1^*$.
        Because of this, it will simply use its own challenge shared key $k^*_1$ if $c^*_1$ is in the decapsulation query by $\adversary{A}$.

        The strategy adopted by $\adversary{B}$ correctly interpolates between games $G_0$ and $G_1$, except when $c^*_1$ would decapsulate to something other than $k_1^*$ in $G_0$.
        We can bound the probability of this inconsistency using the correctness of \kem, which justifies the $\delta$ term in the claim.
        More precisely, we have
        \[
            \abs{\prone{G^{\adversary{A}}_0} - \prone{\G{\kem, \adversary{B}}{0}}} \le \delta,
        \]
        and
        \[
            \prone{G^{\adversary{A}}_1} = \prone{\G{\kem, \adversary{B}}{1}},
        \]
        which means that
        \[
            \abs{\prone{G^{\adversary{A}}_0} - \prone{G^{\adversary{A}}_1}} \leq \advantage{\cca, \adversary{B}}{\kem} + \delta,
        \]
        and $\adversary{B}$ will clearly perform at most as many decapsulation queries to its \kem decapsulation oracle as $\adversary{A}$ makes decapsulation queries.
    \end{itemize}

    \begin{itemize}
    \item[\textbf{Game 2}]
        In this game, all \hqb shared keys that are provided to the adversary computed using $k_1^*$ are replaced with random values.
    \item[\emph{Claim 2}]
        This change should not be noticeable by the adversary $\adversary{A}$.
        If it is, then we can construct an efficient adversary $\adversary{D}$ against the PRF security of $H$ with roughly the same running time as $\adversary{A}$, such that

        \[
            \abs{\prone{G^{\adversary{A}}_1} - \prone{G^{\adversary{A}}_2}} \leq \advantage{\prf, \adversary{D}}{H}.
        \]

    \item[\emph{Proof}]
        This is a simple reduction shown in Figure \ref{fig:adversary-C}.
        This adversary can generate all cryptographic parameters except $k_1^*$, for which it uses its PRF oracle.
        As we can see, when the challenge PRF key is real, then $\adversary{D}$ will be using a real output of the pseudorandom function and hence running $G_1$; whereas if the challenge key is random, $\adversary{D}$ will be using a uniformly sampled output and running $G_2$.
        Hence,

        \[
            \prone{G^{\adversary{A}}_1} = \prone{\G{H, \adversary{D}}{0}},
        \]

        and

        \[
            \prone{G^{\adversary{A}}_2} = \prone{\G{H, \adversary{D}}{1}},
        \]

        which means that

        \[
            \abs{\prone{G^{\adversary{A}}_1} - \prone{G^{\adversary{A}}_2}} \leq \advantage{\prf, \adversary{D}}{H},
        \]

        and $\adversary{D}$ will perform one evaluation query for generating the challenge key and perform at most one evaluation query per decapsulation query by $\adversary{A}$.
    \end{itemize}

    \begin{itemize}
    \item[\textbf{Game 3}]
        We undo the modification introduced in the previous game, but only for keys output by the decapsulation oracle.
    \item[\emph{Claim 3}]
        We justify this hop with a reduction to the PRF property of $H$ very similar to the previous one.

        \[
            \abs{\prone{G^{\adversary{A}}_2} - \prone{G^{\adversary{A}}_3}} \leq \advantage{\prf, \adversary{E}}{H}.
        \]

    \item[\emph{Proof}]
        We construct $\adversary{E}$ as in Figure \ref{fig:adversary-C}, and the analysis is similar to the previous hop.
        \[ \prone{G^{\adversary{A}}_2} = \prone{\prf^{0}_{H, \adversary{E}}}, \]
        and
        \[ \prone{G^{\adversary{A}}_3} = \prone{\prf^1_{H, \adversary{E}}}, \]
        which means that
        \[
            \abs{\prone{G^{\adversary{A}}_2} - \prone{G^{\adversary{A}}_3}} \leq \advantage{\prf, \adversary{E}}{H},
        \]
        and $\adversary{E}$ will perform at most one evaluation query per decapsulation query by $\adversary{A}$ as it will never call the evaluation oracle in its main algorithm.
    \end{itemize}

    \begin{itemize}
    \item[\textbf{Game 4}]
        Finally, we revert the changes introduced in Game 1 and use a real
        key for $k_1^*$ once more.
        This makes the decapsulation oracle identical to Game 0, and the only remaining change in the main experiment is that $k^*$ is random.
    \item[\emph{Claim 4}]
        This hop is justified in a way that is very similar to the jump to Game 1.
        We introduce adversary $\adversary{C}$ against the IND-CCA security of \kem, with roughly the same running time as $\adversary{A}$, in Figure \ref{fig:adversary-B}.
        We claim that,

        \[
            \abs{\prone{G^{\adversary{A}}_3} - \prone{G^{\adversary{A}}_4}} \leq \advantage{\cca, \adversary{C}}{\kem} + \delta \, .
        \]

    \item[\emph{Proof}]
        Again, the reduction to the IND-CCA security of \kem is perfect, except if $c_2^*$ would not decapsulate to $k_1^*$ in Game 4.
        This event can be bound by the correctness of \kem, and the claim follows.
    \item[\emph{Claim 5}]
        $\prone{G^{\adversary{A}}_4} = \prone{\cca^1_{\hqb, \adversary{A}}}$
    \item[\emph{Proof}]
        This follows from the definition of $\cca^1_{\hqb, \adversary{A}}$.

    \end{itemize}

    The theorem follows from collecting all the terms in the claims.

    \begin{gamefigure}{Adversary $\adversary{B}$ and $\mathcolor{\gfour}{\adversary{C}}$ against the IND-CCA security of \kem. Note that both place at most $q_d$ queries to $\dec_O$, one for each decapsulation query placed by $\adversary{A}$.}{fig:adversary-B}
        \begin{games}{2}
            \begin{gameenv}{\textbf{Adversary} $\adversary{B}^{\dec_O(\cdot)}$/$\mathcolor{\gfour}{\adversary{C}^{\dec_O(\cdot)}}(\pk_1, k^*_1, c^*_1)$}
                \State $\sk_2 \sample \honestexponent$
                \State $\pk_2 \assign \nexp(\ngen, \sk_2)$
                \State $\pk \assign (\pk_1, \pk_2)$
                \State $\sk_e \sample \exponent$
                \State $k_1, c^*_1 \sample \kem.\enc(\pk_1)$
                \State $c^*_2 \assign \nexp(\ngen, \sk_e)$
                \State $k^*_2 \assign \nexp(\pk_2, \sk_e)$
                \State $k^* \assign H(\hashinputchallenges)$ \label{line:b-real-hqb}
                \BeginBox[draw=\gfour]
                \State $k^* \sample \bits^{n}$ \label{line:b-sampled-hqb}
                \Comment{$\adversary{C}$}
                \EndBox
                \State $c^* \assign (c_1^*, c_2^*)$
                \State $b' \sample \adversary{A}^{\dec(\cdot)}(\pk, c^*, k^*)$
                \State \Return $b'$
            \end{gameenv}
            \gamebreak
            \begin{gameenv}{\oracle $\dec(c)$}
                \State $\ifthen{c = c^*}$
                \State $\quad$ \Return $\bot$
                \State $(c_1, c_2) \assign c$
                \State $k_2 \assign \nexp(c_2, \sk_2)$
                \State $\ifthen{c_1 \compare c^*_1}$
                \State $\quad k_1 \assign k^*_1$
                \State $\elsedo{k_1 \assign \dec_O(c_1)}$
                \State $\ifthen{k_1 \compare \bot}$
                \State \quad \Return $\bot$
                \State $k \assign H(\hashinput)$
                \State \Return $k$
            \end{gameenv}
        \end{games}
    \end{gamefigure}


    \begin{gamefigure}{Adversary $\adversary{D}$ and $\mathcolor{\gthree}{\adversary{E}}$ against the PRF security of $H$.}{fig:adversary-C}
        \begin{games}{2}
            \begin{gameenv}{\textbf{Adversary} $\adversary{D}^{\Eval_H(\cdot)}$/$\mathcolor{\gthree}{\adversary{E}^{\Eval_H(\cdot)}}()$}
                \State $\sk_1, \pk_1 \sample \kem.\keygen()$
                \State $\sk_2 \sample \honestexponent$
                \State $\pk_2 \assign \nexp(\ngen, \sk_2)$
                \State $\pk \assign (\pk_1, \pk_2)$
                \State $\sk \assign (\sk_1, \sk_2, \pk_2)$
                \State $k_1, c^*_1 \sample \kem.\enc(\pk_1)$
                \State $\sk_e \sample \exponent$
                \State $c^*_2 \assign \nexp(\ngen, \sk_e)$
                \State $k^*_2 \assign \nexp(\pk_2, \sk_e)$
                \State $k^* \assign \Eval_{H}(\text{label} \concat k^*_2 \concat c^*_2\concat \pk_2)$
                \BeginBox[draw=\gthree]
                \State $k^* \sample \bits^{n}$ \label{line:c-hqb-sampled}
                \Comment{$\adversary{E}$}
                \EndBox
                \State $c^* \assign (c_1^*, c_2^*)$
                \State $b' \sample \adversary{A}^{\dec(\cdot)}(\pk, c^*, k^*)$
                \State \Return $b'$
            \end{gameenv}
            \gamebreak
            \begin{gameenv}{\oracle $\dec(c)$}
                \State $\ifthen{c = c^*}$
                \State $\quad$ \Return $\bot$
                \State $(c_1, c_2) \assign c$
                \State $k_2 \assign \nexp(c_2, \sk_2)$
                \State $\ifthen{c_1 \compare c^*_1}$
                \State $\quad k \assign \Eval_{H}(\text{label} \concat k_2 \concat c_2 \concat \pk_2)$
                \State $\quad$ \Return $k$
                \State $k_1 \assign \kem.\dec(\sk_1,c_1)$
                \State $\ifthen{k_1 \compare \bot}$
                \State \quad \Return $\bot$
                \State $k \assign H(\hashinput)$
                \State \Return $k$
            \end{gameenv}
        \end{games}
    \end{gamefigure}
\end{proof}
